\section{Model Convergence and the Limits of Self-Arbitration}

Every empirical and comparative claim in this section is time-sensitive, benchmark-specific, and should be treated as illustrative rather than as a current, load-bearing fact. Model capability comparisons date within months; specific benchmark numbers depend on exact model versions, evaluation harnesses, and scoring methodology that shift constantly. Anywhere this section states a specific comparison without naming the exact models, versions, benchmark, date, and evaluation method involved, treat it as a directional illustration to be independently re-verified before being relied upon, not as settled fact.

The general direction, stated cautiously: performance gaps between flagship models from different labs have narrowed on a number of saturated academic benchmarks, and some open-weight models have closed much of the gap with proprietary flagships on specific coding and reasoning benchmarks. The reasonable directional conclusion is that raw model quality functions increasingly as a converging, commoditizing input rather than a durable point of differentiation on its own --- though the precise magnitude of convergence, and how long any specific ranking holds, changes fast enough that a specific claim (a stated percentage gap, a specific cost multiple) should not be quoted from this report without checking it against current, dated sources first.

Differentiation between products has plausibly moved to three layers, though this too is a framework proposition rather than an established empirical law: native agentic behavior (tool-use reliability, long-horizon planning --- still converging, though with narrow and fast-rotating pockets of advantage); platform and harness philosophy (different labs making different architectural bets, even as underlying protocols for tool connectivity converge toward shared standards); and organizational integration (an organization's own tool graph, evaluation harnesses, and governance layer) --- the layer plausibly hardest to purchase directly from any model vendor, and therefore where durable differentiation is most likely to actually reside.

On the bespoke small-model question: the Terminus-4B case study concerns \textbf{agentic terminal execution} --- a small, specifically post-trained model (built on a 4-billion-parameter base) acting as a subagent that runs terminal commands and interprets their output on behalf of a larger coding agent. The reported finding is that this subagent reduced the token usage of the main agent by roughly 30\% on coding benchmarks with no measured drop in overall agent performance, and that the main agent needed to redo or re-verify the subagent's work less often --- a reasonable, if narrow, proxy for a lower verification tax (Section 7).

The general shape of that finding is worth keeping, stated cautiously. Narrow, purpose-built models can perform very well, sometimes matching or exceeding general frontier models, on closed-form, well-defined, high-volume tasks --- classification against a fixed taxonomy, structured extraction, narrow execution subtasks. They can also underperform general frontier models specifically on tasks requiring open-form or cross-domain judgment, where narrow tuning appears to trade away exactly the broad knowledge that turns out to matter. Fine-tuning \textbf{can} produce an in-distribution/out-of-distribution trade-off, through mechanisms like catastrophic forgetting or overconfidence outside the trained distribution, but the trade-off is not inevitable, and depends heavily on the training data, the specific method used, and how the model is evaluated afterward.

A closed taxonomy does not give a bespoke model a checkable competence boundary for free. Open-set recognition --- correctly identifying that a given input doesn't belong to any of the trained categories at all --- remains a genuinely hard, actively researched problem. Ambiguous cases, drift in the taxonomy itself over time, and setting a sensible abstention threshold are all substantive engineering problems in their own right. A closed label set makes \emph{formal output validation} easier (the output is always one of a known set of labels, which is at least checkable), but it does not automatically tell you whether a given input was actually within the model's competence --- that still requires dedicated testing designed for exactly this question.

The genuinely open question, unaddressed by current evidence either way: how do these dynamics change when the "narrow domain" in question is organization-specific knowledge layered onto a frontier model via retrieval, rather than via full fine-tuning on a narrow training corpus? This retrieval-augmented pattern is the more realistic architecture for most professional-services use cases, and current evidence doesn't clearly distinguish its behavior from either the pure-fine-tuning or pure-frontier-generalist cases already discussed.
